\documentclass[aspectratio=169,10pt]{beamer}\usetheme{metropolis}
\title{Q02: Is Next-Token Generation Intelligence?}\subtitle{Turn a label into testable capability claims}\author{DLS 2026}\date{October 2026}
\begin{document}\maketitle
\begin{frame}{The student question}\begin{block}{Q02}Is next-token generation intelligence? How did intelligence emerge from deep learning?\end{block}\begin{itemize}\item Separate a philosophical label from observable behavior.\item Ask what an experiment would have to show.\end{itemize}\end{frame}
\begin{frame}{One word, many possible claims}\begin{itemize}\item Task performance\item Flexible generalization\item Planning and learning\item Social understanding or consciousness\end{itemize}\alert{Do not let a single word hide which claim is being made.}\end{frame}
\begin{frame}{The objective does not settle the label}\[\min_\theta -\sum_t\log p_\theta(x_{t+1}\mid x_{\leq t})\]\begin{itemize}\item A next-token objective can yield useful task behavior.\item It neither proves nor disproves every sense of ``intelligence.''\item Architecture, data, post-training, prompting, and tools matter too.\end{itemize}\end{frame}
\begin{frame}{Make a claim testable}\textbf{Too broad:} ``The model reasons.''\par\bigskip\textbf{Testable:} On held-out, adversarially paraphrased problems, it achieves a stated score and its verifier-checked artifacts satisfy a stated property.\par\bigskip\alert{Specify: task distribution, metric, baselines, and failure analysis.}\end{frame}
\begin{frame}{What scale results can show}\begin{itemize}\item Performance can improve with data, compute, and capacity.\item Some benchmark curves look threshold-like.\item A threshold may be algorithmic, metric-induced, or both.\end{itemize}\vfill\small Wei et al. (2022); Schaeffer et al. (2023)\end{frame}
\begin{frame}{Three bad inferences}\begin{enumerate}\item Fluency $\Rightarrow$ understanding\item One benchmark $\Rightarrow$ general intelligence\item One failure $\Rightarrow$ no capability\end{enumerate}\vfill Each shortcut discards the task, distribution, or alternative explanation.\end{frame}
\begin{frame}{Mini experiment}\begin{itemize}\item Show a simple rule through examples.\item Hold out combinations that distinguish a rule from memorization.\item Compare a rule-based solver with a lookup baseline.\item Report exact successes and failures.\end{itemize}\alert{The exercise tests a narrow capability, not intelligence as a whole.}\end{frame}
\begin{frame}{Takeaway}\begin{block}{Scientific posture}Treat ``intelligence'' as a family of claims. State the behavior precisely, test it under meaningful distribution shift, and keep conceptual questions separate from measurement claims.\end{block}\vfill\small References: Turing (1950); Brown et al. (2020); Wei et al. (2022); Schaeffer et al. (2023); Bender et al. (2021).\end{frame}
\end{document}
